\section{An expanding parity system}
\label{sec:parity}

We now construct a function that is nonnegative on a product of Boolean
blocks, but is separated from sums of squares of low block degree by a
linear functional.  The blocks record incidence bits of a graph.  At each
vertex these bits have a prescribed parity, and the total prescribed parity
is odd.  Consequently some edge must have different bits at its two ends.
The separating functional nevertheless assigns value zero to the number of
such disagreements.  The construction is a block-coordinate version of the
low-degree functionals for Tseitin contradictions; compare
Grigoriev~\cite{grigoriev}.  We give the construction and its positivity
proof in full.

Fix the constants
\begin{equation}
  d=100,\qquad \beta=\frac1{100},\qquad
  \tau=\frac{\beta}{10}=\frac1{1000}.
  \label{eq:parity-constants}
\end{equation}
All edges of the multigraphs below are distinguished, even when they have
the same endpoints.  Cuts and degrees count these edges with their
multiplicities.  For a multigraph \(H\) and a set \(W\) of its vertices,
\(\delta_H(W)\) consists of the edges with exactly one endpoint in \(W\).
Expansion will select a unique small side whenever a cut has few edges;
this will specify the signs in our functional.

\begin{lemma}[Expanding multigraphs]
  \label{lem:expander}
  For every sufficiently large even integer \(t\), there is a bipartite
  \(d\)-regular multigraph \(H\) on \([t]\) such that
  \begin{equation}
    |\delta_H(W)|\geq \beta d\min\{|W|,t-|W|\}
    \qquad (W\subseteq[t]).
    \label{eq:parity-expansion}
  \end{equation}
  In particular, \(H\) is connected.
\end{lemma}

\begin{proof}
  Write \(t=2l\), split the vertices into two sets of size \(l\), and take
  the union of \(d\) independent uniform perfect matchings between the two
  sides.  Equivalently, choose \(d\) independent uniform permutations of
  \([l]\).

  It suffices to check vertex sets of size \(w\leq l\).  Fix such a
  nonempty set \(W\), containing \(a\) vertices on the first side and
  \(b\) on the second, so that \(a+b=w\).  Always
  \(|\delta_H(W)|\geq d|a-b|\).  Thus a failure of
  \eqref{eq:parity-expansion} is possible only when
  \(|a-b|<.1w\), in which case
  \[
    .45w\leq a,b\leq .55w.
  \]
  If the cut has fewer than \(\beta dw\) edges, fewer than \(.05da\)
  of the \(da\) images of the first-side vertices can lie outside the
  second-side subset: indeed their number is at most
  \(\beta dw\leq (.01/.45)da<.05da\).

  There are at most \(\exp(.21da)\) choices for these exceptional images,
  where an image is specified by its domain vertex and permutation label.
  Here we used
  \[
    \sum_{j\leq .05N}\binom Nj
      \leq \exp\bigl(N[-.05\log(.05)-.95\log(.95)]\bigr)
      \leq \exp(.21N).
  \]
  For any prescribed set of nonexceptional images, sampling without
  replacement within each permutation bounds the probability that they
  all lie in the second-side subset by \((b/l)\) to their number.
  Independence between permutations therefore gives
  \[
    \Pr\bigl[|\delta_H(W)|<\beta dw\bigr]
      \leq \exp(.21da)(.55w/l)^{.95da}.
  \]
  The exponent \(.21+.95\log(.55w/l)\) is negative.  Using
  \(da\geq45w\) and \(\binom{2l}{w}\leq(2el/w)^w\), a union bound
  over all sets of size \(w\) is consequently at most
  \[
    \exp\!\left(
      w\left[\log(2e)+45\bigl(.21+.95\log(.55)\bigr)
                 +41.75\log(w/l)\right]\right)
    \leq e^{-8w}.
  \]
  The last inequality uses \(w/l\leq1\) and
  \(\log(2e)+45(.21+.95\log(.55))<-8\); all logarithms here are
  natural.  Since \(\sum_{w\geq1}e^{-8w}<1\), some choice of the
  permutations satisfies \eqref{eq:parity-expansion}.
  A nontrivial connected component would have an empty cut, so the
  expansion inequality also gives connectivity.
\end{proof}

Fix a graph \(H\) as in Lemma~\ref{lem:expander}, and order the \(d\)
incidences at each vertex.  Choose bits \(b_1,\ldots,b_t\) with odd sum;
one may take \(b_1=1\) and \(b_i=0\) for \(i>1\).  At vertex \(i\),
let
\begin{equation}
  Y_i=\left\{y_i\in\{0,1\}^d:
               \sum_{j=1}^d y_{i,j}=b_i\pmod2\right\},
  \qquad Y=\prod_{i=1}^tY_i.
  \label{eq:parity-domain}
\end{equation}
The first \(d-1\) bits of \(y_i\) are free and determine the last one.
Thus each \(Y_i\) is an affine Boolean cube of size
\(L=2^{d-1}\), equipped with uniform probability.

We use the block characters of Section~\ref{sec:fourier} in these free
coordinates.  Explicitly, set \(\Lambda=\{0,1\}^{d-1}\), with
coordinatewise addition modulo two and zero element \(0\).  For
\(p=(p_1,\ldots,p_t)\in\Lambda^t\), define
\[
  \chi_p(y)=(-1)^{\sum_{i=1}^t\sum_{j=1}^{d-1}p_{i,j}y_{i,j}},
  \qquad |p|_{\mathrm{blk}}=|\{i:p_i\ne0\}|.
\]
These characters form an orthonormal basis for real functions on \(Y\).
A function has \emph{block degree at most \(D\)} if its Fourier
coefficients vanish for \(|p|_{\mathrm{blk}}>D\).

For an edge \(e=uv\), denote its two incidence bits by \(y_{u,e}\) and
\(y_{v,e}\).  Given \(\epsilon>0\), define
\begin{equation}
  f_\epsilon(y)=
    \sum_{e=uv\in E(H)}\mathbf1_{\{y_{u,e}\ne y_{v,e}\}}-\epsilon.
  \label{eq:parity-function}
\end{equation}
The sum of the disagreement indicators is odd.  Indeed, modulo two it
equals the sum of all incidence bits, and
\(\sum_i b_i\equiv1\pmod2\).
In particular, \(f_\epsilon\geq1-\epsilon>0\) when \(0<\epsilon<1\).

\begin{proposition}[A functional separating low-degree squares]
  \label{prop:parity-functional}
  For every sufficiently large even \(t\), let \(H\), \(Y\), and
  \(f_\epsilon\) be as above, and put \(D=\lfloor\tau t\rfloor\).
  There is a real linear functional \(\mathcal D\) on the space of all
  real functions on \(Y\) such that
  \begin{align}
    \mathcal D(1)&=1,
      \label{eq:parity-functional-one}\\
    \mathcal D(g^2)&\geq0
      &&\text{if \(g\) has block degree at most \(D\)},
      \label{eq:parity-functional-positivity}\\
    \mathcal D(f_\epsilon)&=-\epsilon
      &&(\epsilon>0),
      \label{eq:parity-functional-value}\\
    |\mathcal D(h)|&\leq L^t\norm{h}_\infty
      &&\text{for every real function \(h\) on \(Y\)}.
      \label{eq:parity-functional-norm}
  \end{align}
\end{proposition}

\begin{proof}
  We specify \(\mathcal D\) on characters.  First map any set of
  incidences to the edges having an odd number of incidences in that set.
  A character index \(p\) selects the incidences \((i,j)\) for which
  \(j<d\) and \(p_{i,j}=1\); denote their image by
  \(E(p)\subseteq E(H)\).  This incidence map is linear over
  \(\mathbb F_2\), so
  \[
    E(p+p')=E(p)\mathbin\triangle E(p'),
  \]
  where \(\triangle\) denotes symmetric difference.

  If incidence signs agreed across every edge, \(\chi_p\) would be the
  product of the edge signs on \(E(p)\).  Multiplying the vertex parity
  identities over a set \(S\) would prescribe the sign
  \((-1)^{\sum_{i\in S}b_i}\) for the product on \(\delta_H(S)\):
  the signs of edges internal to \(S\) occur twice and cancel.
  Complementary sides of a cut prescribe opposite signs because the
  total prescribed parity is odd.  When \(E(p)\) is a cut and \(p\) has
  low block degree, expansion selects a unique small side and hence a
  value for the functional.

  \emph{Choosing the values on characters.}
  Suppose that \(|p|_{\mathrm{blk}}\leq2D\) and that \(E(p)\) is a cut
  of \(H\).  Since \(|E(p)|\leq2Dd\), expansion gives a side \(S(p)\)
  of this cut with
  \begin{equation}
    |S(p)|\leq\frac{2D}{\beta}
       \leq\frac{2\tau t}{\beta}=\frac t5.
    \label{eq:parity-small-side}
  \end{equation}
  This side is unique: connectivity implies that two vertex sets with
  the same cut are equal or complementary.  Define
  \begin{equation}
    \mathcal D(\chi_p)=
    \begin{cases}
      (-1)^{\sum_{i\in S(p)}b_i},
        & |p|_{\mathrm{blk}}\leq2D
          \text{ and }E(p)\text{ is a cut},\\
      0,&\text{otherwise}.
    \end{cases}
    \label{eq:parity-functional-definition}
  \end{equation}
  Extending linearly defines \(\mathcal D\) on every real function on
  \(Y\).  The trivial character has \(S(0)=\varnothing\), so
  \(\mathcal D(1)=1\).

  \emph{Positivity on low-degree squares.}
  Let \(\mathcal I_D=\{p\in\Lambda^t:|p|_{\mathrm{blk}}\leq D\}\),
  and consider the matrix
  \[
    M_{p,p'}=\mathcal D(\chi_p\chi_{p'})
             =\mathcal D(\chi_{p+p'})
    \qquad(p,p'\in\mathcal I_D).
  \]
  The cut sets of \(H\) form a vector subspace of
  \(\mathbb F_2^{E(H)}\).  Partition \(\mathcal I_D\) according to the
  coset of this subspace containing \(E(p)\).  Entries between different
  classes are zero.  Within a class every entry is a sign, and every
  diagonal entry is one.

  To show that such a block is positive semidefinite, take any three
  indices \(p_1,p_2,p_3\) in it.  Each pair sum has block degree at most
  \(2D\), and so has a small cut side
  \(S_{ab}=S(p_a+p_b)\).  Linearity of both the incidence map and the cut
  map gives
  \[
    \delta_H(S_{12}\triangle S_{23}\triangle S_{31})=\varnothing.
  \]
  The set inside the cut has size at most \(3t/5<t\), by
  \eqref{eq:parity-small-side}.  Connectivity forces it to be empty.
  Therefore
  \[
    M_{p_1,p_2}M_{p_2,p_3}M_{p_3,p_1}=1.
  \]
  Fixing a reference index \(p_0\) in the class and setting
  \(\sigma_p=M_{p,p_0}\), we obtain
  \(M_{p,p'}=\sigma_p\sigma_{p'}\).  Each block is thus a rank-one
  positive semidefinite matrix.  If
  \(g=\sum_{p\in\mathcal I_D}\widehat g(p)\chi_p\), it follows that
  \[
    \mathcal D(g^2)=
      \sum_{p,p'\in\mathcal I_D}\widehat g(p)\widehat g(p')M_{p,p'}
      \geq0.
  \]

  \emph{Values on disagreements.}
  We next show that \(\mathcal D\) assigns value one to the product of
  the two bit signs on every edge.  For a fixed edge \(e=uv\), let
  \(T\subseteq\{u,v\}\) be the endpoints where the incidence of \(e\)
  is in the last position.  On \(Y_i\) the parity constraint gives
  \[
    (-1)^{y_{i,d}}=(-1)^{b_i}
                        \prod_{j=1}^{d-1}(-1)^{y_{i,j}}.
  \]
  Using this identity at the vertices in \(T\), we can write
  \[
    (-1)^{y_{u,e}+y_{v,e}}
       =(-1)^{\sum_{i\in T}b_i}\chi_p(y)
       \qquad\text{with }|p|_{\mathrm{blk}}\leq2.
  \]
  Before these replacements, the two selected incidences belong to the
  same edge and give the empty edge set under the incidence map.
  Replacing a last incidence by the other \(d-1\) incidences toggles the
  full star at that vertex.  Its image under the incidence map is
  \(\delta_H(\{i\})\).  Hence \(E(p)=\delta_H(T)\), including when
  \(u\) and \(v\) have parallel edges.  For sufficiently large \(t\),
  we have \(2\leq2D\) and \(|T|\leq2\leq t/5\), so \(T\) is the
  small side used in \eqref{eq:parity-functional-definition}.
  The two signs cancel, giving
  \[
    \mathcal D\bigl((-1)^{y_{u,e}+y_{v,e}}\bigr)=1.
  \]
  Since
  \(\mathbf1_{\{y_{u,e}\ne y_{v,e}\}}
       =(1-(-1)^{y_{u,e}+y_{v,e}})/2\), every disagreement indicator
  has functional value zero.  This proves
  \(\mathcal D(f_\epsilon)=-\epsilon\).

  \emph{Norm bound.}
  Finally, all \(L^t\) character values in
  \eqref{eq:parity-functional-definition} have absolute value at most
  one.  For any real function \(h\), uniform orthonormal Fourier
  expansion gives
  \(\widehat h(p)=\E_{y\in Y}h(y)\chi_p(y)\), and therefore
  \(|\widehat h(p)|\leq\norm{h}_\infty\).  Thus
  \[
    |\mathcal D(h)|
      \leq\sum_{p\in\Lambda^t}|\widehat h(p)|
      \leq L^t\norm{h}_\infty,
  \]
  completing the proof.
\end{proof}
